% Chunk: \S8.2, Constraints and Gauge Conditions.
% Source: Weinberg QFT Vol. I, physical pp. 366--369 (printed pp. 343--346),
% from the \S8.2 heading through the paragraph immediately before \S8.3.
% Coverage: equations (8.2.1)--(8.2.9); four footnotes; reference [2].

\subsection{Constraints and Gauge Conditions}

There are aspects of electrodynamics that stand in the way of quantizing
the theory as we did for various theories of massive particles in the
previous chapter. As usual, we may define the canonical conjugates to the
electromagnetic vector potential by
\begin{equation}
  \Pi^\mu\equiv\frac{\partial\mathcal L}{\partial(\partial_0A_\mu)}.
  \tag{8.2.1}\label{eq:8.2.1}
\end{equation}
Quantization by the usual rules would give
\[
  [A_\mu(\mathbf x,t),\Pi^\nu(\mathbf y,t)]
  =\ii\delta_\mu{}^\nu\delta^3(\mathbf x-\mathbf y).
\]
But this is not possible here, because $A_\mu$ and $\Pi^\nu$ are subject to
several constraints.

The first constraint arises from the fact that the Lagrangian density is
independent\footnote{For $\mathcal L_\gamma=-F_{\mu\nu}F^{\mu\nu}/4$, we have
$\partial\mathcal L_\gamma/\partial(\partial_0A_\mu)=-F^{0\mu}$, which vanishes
for $\mu=0$ because $F^{\mu\nu}$ is antisymmetric. For matter Lagrangians
$\mathcal L_M$ that involve only $\Psi^\ell$ and $D_\mu\Psi^\ell$, the
prescription \eqref{eq:8.1.21} tells us that $\mathcal L_M$ does not depend on any
derivatives of any $A^\mu$. Even if the matter Lagrangian depends also on
$F_{\mu\nu}$, $\partial\mathcal L_M/\partial(\partial_\mu A_\nu)$ will be again
antisymmetric in $\mu$ and $\nu$, and therefore will vanish for $\mu=\nu=0$.}
of the time-derivative of $A_0$, and therefore
\begin{equation}
  \Pi^0(x)=0.
  \tag{8.2.2}\label{eq:8.2.2}
\end{equation}
This is called a \textit{primary constraint}, because it follows directly from
the structure of the Lagrangian. There is also a \textit{secondary constraint}
here, which follows from the field equation for the quantity fixed by the
primary constraint:\footnote{As usual, $i,j$, etc. run over the values $1,2,3$.}
\begin{equation}
  \partial_i\Pi^i=-\partial_i\frac{\partial\mathcal L}{\partial F_{i0}}
  =-\frac{\partial\mathcal L}{\partial A_0}\equiv-J^0,
  \tag{8.2.3}\label{eq:8.2.3}
\end{equation}
the time-derivative term dropping out because $F_{00}=0$. Even though
the matter Lagrangian may generally depend on $A^0$, the charge density
depends only on the canonical matter fields\footnote{Due to exhaustion of
alphabetic resources, I have had to adopt a notation here that is different
from that of the previous chapter. The symbols $Q^n$ and $P_n$ are now
reserved for the canonical matter fields and their canonical conjugates,
respectively, while the canonical electromagnetic fields and canonical
conjugates are $A_i$ and $\Pi_i$.} $Q^n$ and their canonical conjugates $P_n$:
\begin{equation}
  J^0=-\ii\sum_\ell
  \frac{\partial\mathcal L}{\partial(\partial_0\Psi^\ell)}q_\ell\Psi^\ell
  =-\ii\sum_n P_nq_nQ^n.
  \tag{8.2.4}\label{eq:8.2.4}
\end{equation}
Hence Eq.~\eqref{eq:8.2.3} is a functional relation among canonical variables.
Both Eq.~\eqref{eq:8.2.2} and Eq.~\eqref{eq:8.2.3} are inconsistent with the usual
assumptions that
$[A_\mu(\mathbf x,t),\Pi^\nu(\mathbf y,t)]
=\ii\delta_\mu{}^\nu\delta^3(\mathbf x-\mathbf y)$ and
$[Q^n(\mathbf x,t),\Pi^\nu(\mathbf y,t)]
=[P_n(\mathbf x,t),\Pi^\nu(\mathbf y,t)]=0$.

We encountered a similar problem in the theory of the massive vector
field. In that case we found two equivalent ways of dealing with it:
either by the method of Dirac brackets or, more directly, by treating only
$A_i$ and $\Pi^i$ as canonical variables, solving the analog of Eq.~\eqref{eq:8.2.3} to
calculate $A^0$ in terms of these variables. It is clear that here we cannot
use Dirac brackets; the constraint functions $\chi$ here are $\Pi^0$ and
$\partial_i\Pi_i-J^0$ (as compared with $\partial_i\Pi_i-m^2A^0-J^0$) and these
obviously have vanishing Poisson brackets. In Dirac's terminology, the
constraints \eqref{eq:8.2.2} and \eqref{eq:8.2.3} are \textit{first class}. Nor can we eliminate
$A^0$ as a dynamical variable by solving for it in terms of the other variables.
Instead of giving $A^0$ for all time, Eq.~\eqref{eq:8.2.3} is a mere initial condition;
if Eq.~\eqref{eq:8.2.3} is satisfied at one time, then it is satisfied for all times,
because (using the field equations for the other fields $A^i$), we have
\begin{align*}
  \partial_0\left[\partial_i\frac{\partial\mathcal L}{\partial F_{i0}}-J^0\right]
  &=-\partial_i\partial_0\frac{\partial\mathcal L}{\partial F_{0\ii}}
  -\partial_0J^0\\
  &=+\partial_i\partial_j\frac{\partial\mathcal L}{\partial F_{ji}}
  -\partial_iJ^i-\partial_0J^0,
\end{align*}
and the current conservation condition then gives
\begin{equation}
  \partial_0\left[\partial_i\frac{\partial\mathcal L}{\partial F_{i0}}-J^0\right]=0.
  \tag{8.2.5}\label{eq:8.2.5}
\end{equation}
It should not be surprising that we still have four components of $A^\mu$
with only three field equations, because this theory has a local gauge
symmetry that makes it, in principle, impossible to infer the values of the
fields at arbitrary times from their values and rates of change at any one
time. Given any solution $A_\mu(\mathbf x,t)$ of the field equations, we can
always find another solution $A_\mu(\mathbf x,t)+\partial_\mu\epsilon(\mathbf x,t)$
with the same value and time-derivative at $t=0$ (by choosing $\epsilon$ so that
its first and second derivatives vanish there) but which differs from
$A_\mu(\mathbf x,t)$ at later times.

Because of this partial arbitrariness of $A_\mu(\mathbf x,t)$, it is not possible
to apply the canonical quantization procedure directly to $A_\mu$ (or, as for
finite mass, to $\mathbf A$). Of the various approaches to this difficulty, two
are particularly useful. One is the modern method of gauge-invariant
quantization, to be discussed in Volume II. The other, which will be followed
here, is to exploit the gauge invariance of the theory, to `choose a gauge.'
That is, we make a finite gauge transformation
\[
  A_\mu(x)\longrightarrow A_\mu(x)+\partial_\mu\lambda(x),
  \qquad
  \Psi_\ell(x)\longrightarrow\exp\!\bigl(\ii q_\ell\lambda(x)\bigr)\Psi_\ell(x)
\]
to impose a condition on $A_\mu(x)$ that will allow us to apply the methods
of canonical quantization. There are various gauges that have been found
useful in various applications:\footnote{Here $\Phi$ is any complex scalar
field with $q\ne0$; this gauge condition is used when the gauge symmetry is
spontaneously broken by a non-vanishing vacuum expectation value of $\Phi$.}
\begin{align*}
  \text{Lorentz (or Landau) gauge:}\quad&\partial_\mu A^\mu=0,\\
  \text{Coulomb gauge:}\quad&\boldsymbol\nabla\cdot\mathbf A=0,\\
  \text{Temporal gauge:}\quad&A^0=0,\\
  \text{Axial gauge:}\quad&A^3=0,\\
  \text{Unitarity gauge:}\quad&\Phi\ \text{real}.
\end{align*}
The canonical quantization procedure works most easily in the axial or
Coulomb gauge, but of course Coulomb gauge keeps manifest rotation
invariance in a way that axial gauge does not, so we will adopt Coulomb
gauge here~\hyperref[ref:8.2]{[2]}.

To check that this is possible, note that if $A^\mu$ does not satisfy the
Coulomb gauge condition, then the gauge-transformed field
$A^\mu+\partial^\mu\lambda$ will, provided we choose $\lambda$ so that
$\nabla^2\lambda=-\boldsymbol\nabla\cdot\mathbf A$. From now on, we assume
that this transformation has been made, so that
\begin{equation}
  \boldsymbol\nabla\cdot\mathbf A=0.
  \tag{8.2.6}\label{eq:8.2.6}
\end{equation}

It will be convenient henceforth to limit ourselves to theories in which
the matter Lagrangian $\mathcal L_M$ may depend on matter fields and their
time-derivatives and also on $A^\mu$ but not on derivatives of $A^\mu$. (The
standard theories of the electrodynamics of scalar and Dirac fields have
Lagrangians of this type.) Then the only term in the Lagrangian that depends
on $F_{\mu\nu}$ is the kinematic term $-\tfrac14F_{\mu\nu}F^{\mu\nu}$, and the
constraint equation~\eqref{eq:8.2.3} reads
\begin{equation}
  -\partial_iF^{i0}=J^0.
  \tag{8.2.7}\label{eq:8.2.7}
\end{equation}
Together with the Coulomb gauge condition \eqref{eq:8.2.6}, this yields
\begin{equation}
  -\nabla^2A^0=J^0,
  \tag{8.2.8}\label{eq:8.2.8}
\end{equation}
which can be solved to give
\begin{equation}
  A^0(\mathbf x,t)=\int \dd^3y\,
  \frac{J^0(\mathbf y,t)}{4\pi\lvert\mathbf x-\mathbf y\rvert}.
  \tag{8.2.9}\label{eq:8.2.9}
\end{equation}
The remaining degrees of freedom are $A^i$, with $i=1,2,3$, subject to the
gauge condition $\boldsymbol\nabla\cdot\mathbf A=0$.

As mentioned earlier, the charge density depends only on the canonical
matter fields $Q^n$ and their canonical conjugates $P_n$, so Eq.~\eqref{eq:8.2.9}
represents an explicit solution for the auxiliary field $A^0$.
